% Macro-heavy thesis preamble: the case that breaks naive block parsers.
\documentclass[12pt,twoside,openright]{report}
\usepackage{amsmath,amssymb,mathtools}
\usepackage{xparse}
\usepackage{etoolbox}

% --- Notation shorthands --------------------------------------------------
\newcommand{\R}{\mathbb{R}}
\newcommand{\N}{\mathbb{N}}
\newcommand{\E}{\mathbb{E}}
\newcommand{\Prob}{\mathbb{P}}
\renewcommand{\vec}[1]{\boldsymbol{#1}}
\providecommand{\abs}[1]{\left\lvert #1 \right\rvert}
\DeclareMathOperator{\argmin}{arg\,min}
\DeclareMathOperator*{\esssup}{ess\,sup}

% A macro with an optional argument, defined the xparse way.
\NewDocumentCommand{\norm}{o m}{%
  \IfNoValueTF{#1}{\left\lVert #2 \right\rVert}{\left\lVert #2 \right\rVert_{#1}}%
}

\newenvironment{remark}[1][Remark]{%
  \par\medskip\noindent\textbf{#1.}\enspace%
}{%
  \par\medskip%
}

\newcommand{\citefig}[2]{Figure~\ref{#1} (see also \cite{#2})}

\begin{document}

\chapter{Preliminaries}
\label{ch:prelim}

We work throughout on a probability space $(\Omega, \mathcal{F}, \Prob)$ and take
all random elements to be $\R^d$-valued. For $x \in \R^d$ write $\norm{x}$ for the
Euclidean norm and $\norm[1]{x}$ for the $\ell_1$ norm.

\section{Standing assumptions}

\begin{remark}[Convention]
Unless stated otherwise, $\E$ denotes expectation under $\Prob$.
\end{remark}

The estimator is $\hat\theta = \argmin_{\theta} \abs{Q(\theta)}$, and we assume
$\esssup_\omega \norm{X(\omega)} < \infty$.

\input{chapters/background}
\include{chapters/method}

\section{A displayed system}

\begin{align}
  \label{eq:system}
  \dot{x}(t) &= A x(t) + B u(t), \\
  y(t)       &= C x(t).
\end{align}

Note $\citefig{fig:block}{astrom2008}$ for the block diagram.

\end{document}
